% -----------------------------------------------------------
\section{Truth}
% -----------------------------------------------------------
	\label{TRUTH}
	
% % ----------------------
% \subsection{See also}
% % ----------------------

% ----------------------
\subsection{1828 American Dictionary of the English Language} % *1828
% ----------------------
	\label{TRUTH-1828}

	\begin{compactitem}
		\item[1] Conformity to fact or reality; exact accordance with that which is, or has been, or shall be. The \textit{truth} of history constitutes its whole value. We rely on the \textit{truth} of the scriptural prophecies.
		\item[2] True state of facts or things. The duty of a court of justice is to discover the \textit{truth} Witnesses are sworn to declare the \textit{truth} the whole \textit{truth} and nothing but the \textit{truth}.
		\item[3] Conformity of words to thoughts, which is called moral \textit{truth}.
		\item[4] Veracity; purity from falsehood; practice of speaking \textit{truth}; habitual disposition to speak \textit{truth}; as when we say, a man is a man of \textit{truth}.
		\item[5] Correct opinion.
		\item[6] Fidelity; constancy.
		\item[7] Honesty; virtue.
		\item[8] Exactness; conformity to rule.
		\item[9] Real fact of just principle; real state of things. There are innumerable \textit{truth}s with which we are not acquainted.
		\item[10] Sincerity.
		\item[11] The \textit{truth} of God, is his veracity and faithfulness.
		\item[12] Jesus Christ is called the \textit{truth} John 14.
		\item[13] It is sometimes used by way of concession.
	\end{compactitem}

% % ----------------------
% \subsection{Oxford English Dictionary} % *OED
% % ----------------------
	% \label{TRUTH-OED}

	% \begin{compactitem}
		% \item[]
	% \end{compactitem}

% % ----------------------
% \subsection{Scriptural Usage} % *SCRIPTURES
% % ----------------------
	% \label{TRUTH-SCRIPTURES}

	% % -----
	% \subsubsection{Old Testament} % *OT
	% % -----
		% \label{TRUTH-OT}
	
		% % --
		% \paragraph{KJV}
		% % --
			% \label{TRUTH-OT-KJV}
	
			% \begin{tabular}{ll}
				% Total occurrences in the OT & 
			% \end{tabular}
			
			% \begin{compactdesc}
				% \item[]
			% \end{compactdesc}
			
		% % --
		% \paragraph{Hebrew Bible}
		% % --
			% \label{TRUTH-HEBREW}
		
			% %
			% \subparagraph{\texorpdfstring{\he{sample word}}{}} % paste actual hebrew text here
			% %
				% \label{PAGA} % whatever is in step bible without periods
			
				% \begin{compactdesc}
					% \item[HALOT , PDF ] \textbf{} % Use the 1994 PDF
						% \begin{compactitem}
							% \item[1]
						% \end{compactitem}
					% \item[BDB , PDF ] \textbf{}
						% \begin{compactitem}
							% \item[1]
						% \end{compactitem}
					% \item[DCH PDF ] \textbf{}
						% \begin{compactitem}
							% \item[1]
						% \end{compactitem}
				% \end{compactdesc}
				
				% \begin{compactdesc}
					% \item[Some OT uses] \textbf{}
						% \begin{compactdesc}
							% \item[]
						% \end{compactdesc}
				% \end{compactdesc}
			
		% % --
		% \paragraph{Septuagint}
		% % --
			% \label{TRUTH-LXX}
		
			% %
			% \subparagraph{\texorpdfstring{\gr{sample word}}{}}
			% %
		
	% -----
	\subsubsection{New Testament} % *NT
	% -----
		\label{TRUTH-NT}
	
		% --
		\paragraph{KJV}
		% --
			\label{TRUTH-NT-KJV}
			
	
			% \begin{tabular}{ll}
				% Total occurrences in the NT & 
			% \end{tabular}
			
			\begin{compactdesc}
				\item[{\hyperref[JOHN-1-9]{John 1:9}}]
				\item[{\hyperref[JOHN-1-14]{John 1:14}}]
				\item[{\hyperref[JOHN-1-17]{John 1:17}}]
				\item[{\hyperref[JOHN-3-21]{John 3:21}}]
				\item[{\hyperref[JOHN-3-33]{John 3:33}}]
				\item[{\hyperref[JOHN-4-23]{John 4:23--24}}]
				\item[{\hyperref[JOHN-5-31]{John 5:31--33}}]
				\item[{\hyperref[JOHN-6-32]{John 6:32}}]
				\item[{\hyperref[JOHN-7-18]{John 7:18}}]
				\item[{\hyperref[JOHN-7-28]{John 7:28}}]
				\item[{\hyperref[JOHN-8-13]{John 8:13--17}}]
				\item[{\hyperref[JOHN-8-26]{John 8:26}}]
				\item[*{\hyperref[JOHN-8-32]{John 8:32}}]
					\begin{compactitem}
						\item This is the ``truth shall make you free'' verse.
						\item As noted in this list, there is substantial discussion of truth throughout this chapter. It also talks about what is ``true.''
					\end{compactitem}
				\item[{\hyperref[JOHN-8-40]{John 8:40--45}}]
				\item[*{\hyperref[JOHN-14-6]{John 14:6}}]
				\item[{\hyperref[JOHN-14-17]{John 14:17}}]
				\item[{\hyperref[JOHN-15-1]{John 15:1}}]
				\item[{\hyperref[JOHN-15-26]{John 15:26}}]
				\item[{\hyperref[JOHN-16-7]{John 16:7-13}}]
				\item[{\hyperref[JOHN-17-3]{John 17:3}}]
				\item[*{\hyperref[JOHN-17-17]{John 17:17}}]
				\item[*{\hyperref[JOHN-17-19]{John 17:19}}]
			\end{compactdesc}
			
		% % --
		% \paragraph{Greek New Testament} % *GREEK
		% % --
			% \label{TRUTH-GREEK}
		
			% %
			% \subparagraph{\texorpdfstring{\gr{sample word}}{}}
			% %
				% \label{KATALLAGE}
			
				% \begin{compactdesc}
					% \item[BDAG , PDF ] \textbf{}
						% \begin{compactitem}
							% \item 
						% \end{compactitem}
					% \item[CGL , PDF ] \textbf{}
						% \begin{compactitem}
							% \item 
						% \end{compactitem}
					% \item[LSJ , PDF ] \textbf{}
						% \begin{compactitem}
							% \item 
						% \end{compactitem}
				% \end{compactdesc}
				
				% \begin{compactdesc}
					% \item[Some NT uses] \textbf{}
						% \begin{compactdesc}
							% \item[]
						% \end{compactdesc}
				% \end{compactdesc}
		
	% -----
	\subsubsection{Book of Mormon} % *BOM
	% -----
		\label{TRUTH-BOM}
	
		% \begin{tabular}{ll}
			% Total occurrences in the BOM & 
		% \end{tabular}
		
		\begin{compactdesc}
			\item[{\hyperref[3NEPHI-5-18]{3 Nephi 5:18}}]
		\end{compactdesc}
		
	% % -----
	% \subsubsection{Doctrine and Covenants} % *DC
	% % -----
		% \label{TRUTH-DC}
	
		% \begin{tabular}{ll}
			% Total occurrences in the DC & 
		% \end{tabular}
		
		% \begin{compactdesc}
			% \item[]
		% \end{compactdesc}
		
	% % -----
	% \subsubsection{Pearl of Great Price} % *PGP
	% % -----
		% \label{TRUTH-PGP}
	
		% \begin{tabular}{ll}
			% Total occurrences in the PGP & 
		% \end{tabular}
		
		% \begin{compactdesc}
			% \item[]
		% \end{compactdesc}
		
% ----------------------
\subsection{Key Phrases} % *KEYPHRASES *PHRASES
% ----------------------
	\label{TRUTH-PHRASES}

	\begin{compactdesc}
	
		\item[full of grace and truth] \textbf{13} | \hyperref[JOHN-1-14]{John 1:14}, \hyperref[2NEPHI-2-6]{2 Nephi 2:6}; \hyperref[ALMA-5-48]{Alma 5:48}; \hyperref[ALMA-9-26]{Alma 9:26}; \hyperref[ALMA-13-9]{Alma 13:9}; \hyperref[DC66-12]{DC 66:12}; \hyperref[DC84-102]{DC 84:102}; \hyperref[DC93-11]{DC 93:11}; \hyperref[MOSES-1-6]{Moses 1:6}, \hyperref[MOSES-1-32]{32}; \hyperref[MOSES-5-7]{Moses 5:7}; \hyperref[MOSES-6-52]{Moses 6:52}; \hyperref[MOSES-7-11]{Moses 7:11}
			\begin{compactdesc}
				\item[Bible anchor verses] \phantom{.}
					\begin{compactdesc}
						\item[John 1:14]
					\end{compactdesc}
				% \item[Commentary] ---
			\end{compactdesc}
	
		\item[grace and truth] \textbf{15} | \hyperref[JOHN-1-14]{John 1:14}, \hyperref[JOHN-1-17]{17}; \hyperref[2NEPHI-2-6]{2 Nephi 2:6}; \hyperref[ALMA-5-48]{Alma 5:48}; \hyperref[ALMA-9-26]{Alma 9:26}; \hyperref[ALMA-13-9]{Alma 13:9}; \hyperref[DC50-40]{DC 50:40}; \hyperref[DC66-12]{DC 66:12}; \hyperref[DC84-102]{DC 84:102}; \hyperref[DC93-11]{DC 93:11}; \hyperref[MOSES-1-6]{Moses 1:6}, \hyperref[MOSES-1-32]{32}; \hyperref[MOSES-5-7]{Moses 5:7}; \hyperref[MOSES-6-52]{Moses 6:52}; \hyperref[MOSES-7-11]{Moses 7:11}
			\begin{compactdesc}
				\item[Bible anchor verses] \phantom{.}
					\begin{compactdesc}
						\item[John 1:14]
						\item[John 1:17]
					\end{compactdesc}
				\item[Commentary] I wonder if one of the keys to understanding this phrase, besides the uses in John, is the Hebrew phrase \he{.EsEd wE'E:mEt}, and more specifically \he{rab--.hEsEd wE'E:mEt} at the end of \hyperref[EXODUS-34-6]{Exodus 34:6}. I have seen some of the sources gloss that phrase as ``reliability,'' basically taking the \he{.EsEd wE'E:mEt} as a hendiadys. That makes me think of at least two possibilities for ``grace and truth'' (and also ``full of grace and truth''): either take it as hendiadys in a similar way (and maybe even gloss it as ``reliability''); or don't take it as hendiadys and instead expand our understanding of the phrase by looking at OT instances of this Hebrew phrase, along with possible meanings of \he{.hEsEd}. I think I'm inclined to the latter but I would also say to be wary here because you don't want to lose the connections with restoration scripture that are clearly built on the John use of ``grace and truth,'' especially in places like \hyperref[DC93]{DC 93}.
					\begin{compactitem}
						\item See TLOT PDF 605 bottom for more on the Exodus verse.
					\end{compactitem}
			\end{compactdesc}
			
		\item[just and true] \textbf{10} | \hyperref[REVELATIONNT-15-3]{Revelation 15:3}; \hyperref[1NEPHI-14-23]{1 Nephi 14:23}; \hyperref[MOSIAH-2-35]{Mosiah 2:35}; \hyperref[MOSIAH-4-12]{Mosiah 4:12}; \hyperref[ALMA-18-34]{Alma 18:34}; \hyperref[ALMA-29-8]{Alma 29:8}; \hyperref[MORONI-10-6]{Moroni 10:6}; \hyperref[DC20-30]{DC 20:30}, \hyperref[DC20-31]{31}; \hyperref[DC76-53]{DC 76:53}
			\begin{compactdesc}
				\item[Bible anchor verses] \hyperref[REVELATIONNT-15-3]{Revelation 15:3}, \gr{δίκαιαι καὶ ἀληθιναὶ αἱ ὁδοί σου}
				\item[Commentary] The list of ten entries above are places where you find the exact phrase ``just and true.'' If you broaden the search to ``just* tru*'', you get a total of 38 hits, most of which aren't super interesting. But you do get \hyperref[3NEPHI-5-18]{3 Nephi 5:18}, which includes ``a just and a true record,'' and so belongs with the ten entries above.
			\end{compactdesc}
			
		\item[light and truth] \textbf{6} | \hyperref[DC93-36]{DC 93:36}, \hyperref[DC93-37]{37}, \hyperref[DC93-39]{39}, \hyperref[DC93-40]{40}, \hyperref[DC93-42]{42}; \hyperref[DC124-9]{DC 124:9}
			\begin{compactdesc}
				\item[Bible anchor verses] ---
				\item[Commentary] 
			\end{compactdesc}
			
		\item[light of truth] \textbf{2} | \hyperref[DC88-6]{DC 88:6}; \hyperref[DC93-29]{DC 93:29}
			\begin{compactdesc}
				\item[Bible anchor verses] ---
				\item[Commentary] See ``light and truth'' above.
			\end{compactdesc}
	
		\item[spirit of truth] See this phrase at \hyperref[SPIRIT-PHRASES]{Spirit}.
			% \begin{compactdesc}
				% \item[Bible anchor verses] \phantom{.}
					% \begin{compactdesc}
						% \item[Romans 5:5] \gr{}
					% \end{compactdesc}
				% \item[Commentary] ---
			% \end{compactdesc}
			
	\end{compactdesc}
	
% % ----------------------
% \subsection{Bible Dictionaries} % *DICTIONARIES
% % ----------------------
	% \label{TRUTH-BD}

% ----------------------
\subsection{Restoration Usage} % *RESTORATION
% ----------------------
	\label{TRUTH-RESTORATION}

	% % -----
	% \subsubsection{Joseph Smith} % *JS
	% % -----
		% \label{TRUTH-JS}
	
		% \begin{compactdesc}
			% \item[{\hyperref[KEY]{Date/Source}}]
		% \end{compactdesc}
		
	% -----
	\subsubsection{Temple Ordinances}
	% -----
		\label{TRUTH-TEMPLE}
	
		\begin{compactdesc}
			\item[{\hyperref[TEMPLE-ENDOWMENT-ENDOW-PRE-1990-VEIL]{Discussion of garment symbols}}] This one on the left is the mark of the compass. It is placed in the garment over the left breast, suggesting to the mind an undeviating course leading to eternal life; a constant reminder that desires, appetites, and passions are to be kept within the bounds the Lord has set; and that all truth may be circumscribed into one great whole.
		\end{compactdesc}
	
	% % -----
	% \subsubsection{Brigham Young} % *BY
	% % -----
		% \label{TRUTH-BY}
	
		% \begin{compactdesc}
			% \item[{\hyperref[KEY]{Date/Source}}]
		% \end{compactdesc}
	
% % ----------------------
% \subsection{Commentary and Studies} % *COMMENTARY *STUDIES
% % ----------------------
	% \label{TRUTH-COMMENTARY}

	% % -----
	% \subsubsection{Key Points}
	% % -----
		% \label{TRUTH-KEYS}
	
		% \begin{compactitem}
			% \item
		% \end{compactitem}
		
	% -----
	\subsubsection{Truth and goodness}
	% -----
	
		See this study \hyperref[GOODNESS-truth]{here}.